\documentclass[11pt]{article}
\usepackage[margin=1in]{geometry}
\usepackage{amsmath,amssymb,amsthm}
\usepackage{hyperref}

\newtheorem{theorem}{Theorem}
\newtheorem{lemma}[theorem]{Lemma}
\newtheorem{corollary}[theorem]{Corollary}
\newtheorem{proposition}[theorem]{Proposition}

\title{A Disproof of the Fanelli--Moscardelli Stability Threshold Conjecture for Odd Markets}
\author{Salem Levinstein}
\date{November 2026}

\begin{document}
\maketitle

\begin{abstract}
In the approximate stable roommates problem with complete strict preferences on $N$ agents, a pair $\{i,j\}$ $k$-blocks a matching $M$ if both agents improve their preference rank by at least $k$ positions (with being single assigned rank $N$). Let $K(N)$ denote the least positive integer $k$ such that every complete profile on $N$ agents admits a $k$-stable matching. Fanelli and Moscardelli (ECAI 2025) proved $K(N) \le \lceil N/2 \rceil$, verified equality for $N \le 6$, and conjectured that $K(N) = \lceil N/2 \rceil$ for all sufficiently large $N$. We disprove this conjecture by giving an explicit constructive proof that $K(2m+1) \le m < \lceil (2m+1)/2 \rceil$ for every integer $m \ge 3$ (sharp at $m=2$, where $K(5)=3$).
\end{abstract}

\section{Introduction and Definitions}

Let $V$ be a set of $N$ agents. A \emph{complete strict preference profile} assigns to each agent $i \in V$ a bijection $\mathrm{rank}_i : V \to \{1, \dots, N\}$ with $\mathrm{rank}_i(i) = N$. A \emph{matching} $M$ is an involution on $V$; agents with $M(i) = i$ are \emph{single}. Write $\rho_i := \mathrm{rank}_i(M(i))$, so $\rho_i = N$ if $i$ is single and $\rho_i \le N-1$ otherwise.

Following Fanelli and Moscardelli~\cite{FM25}, an unordered pair $\{i,j\}$ ($i \neq j$, $M(i) \neq j$) is a \emph{$k$-blocking pair} for $M$ if
\[
\mathrm{rank}_i(j) + k \le \rho_i \qquad \text{and} \qquad \mathrm{rank}_j(i) + k \le \rho_j.
\]
A matching is \emph{$k$-stable} if it has no $k$-blocking pair, and $K(N)$ is the least integer $k \ge 1$ such that every complete profile on $N$ agents admits a $k$-stable matching.

Fanelli and Moscardelli~\cite[Theorem~1]{FM25} proved $K(N) \le \lceil N/2 \rceil$ for all $N \ge 3$, established an $\Omega(\sqrt{N})$ lower bound~\cite[Theorem~3]{FM25}, verified $K(N) = \lceil N/2 \rceil$ for $3 \le N \le 6$, and conjectured~\cite[Section~5]{FM25} that $K(N) = \lceil N/2 \rceil$ for all sufficiently large $N$. We refute this conjecture for every odd $N \ge 7$.

\begin{theorem}\label{thm:main}
For every integer $m \ge 3$ and $N = 2m + 1$, every complete strict preference profile on $N$ agents admits an $m$-stable matching. Consequently,
\[
K(2m+1) \le m < m + 1 = \lceil (2m+1)/2 \rceil \qquad \text{for all } m \ge 3.
\]
\end{theorem}

The hypothesis $m \ge 3$ is sharp, since $K(5) = 3$, and is used only in Section~\ref{sec:allk2} to ensure $|\mathrm{out}(z)| = m - 2 \ge 1$. Accompanying DRAT-checked SAT computations and lower-bound witnesses (available in the companion repository) further establish the exact values $K(7)=3$, $K(8)=3$, $K(9)=4$, and $K(11)=5$.

\paragraph{Notation.} Fix $m \ge 3$ and $N = 2m+1$. For each $v \in V$, set $T_v := \{u \in V : \mathrm{rank}_v(u) \le m\}$ (so $|T_v| = m$ and $v \notin T_v$). Let $G = (V, E(G))$ be the \emph{mutual-top-$m$ graph}, with $uv \in E(G)$ iff $u \in T_v$ and $v \in T_u$, and let $N(v)$ denote the neighborhood of $v$ in $G$. Let $x_v \in V$ be $v$'s $(m+1)$-th choice ($\mathrm{rank}_v(x_v) = m+1$) and $f(v) \in V$ be $v$'s first choice ($\mathrm{rank}_v(f(v)) = 1$). Given a matching $M$ with a single unmatched agent $u$, a matched agent $i \neq u$ is \emph{content w.r.t.\ $j$} if $\rho_i \le \mathrm{rank}_i(j) + m - 1$. If $\rho_i \le m$ (in particular, if $i$ is matched along a $G$-edge), then $i$ is content w.r.t.\ every agent.

\begin{lemma}[Stability criterion]\label{lem:crit}
A matching $M$ on $N = 2m+1$ agents with exactly one single agent $u$ is $m$-stable if and only if:
\begin{enumerate}
  \item[(a)] every $G$-edge $\{i,j\} \notin M$ has an endpoint $i \neq u$ that is content w.r.t.\ $j$; and
  \item[(b)] $\rho_{x_u} \le \mathrm{rank}_{x_u}(u) + m - 1$ (which holds automatically if $\mathrm{rank}_{x_u}(u) \ge m+1$).
\end{enumerate}
\end{lemma}

\begin{proof}
If both $i,j$ are matched, then $\rho_i, \rho_j \le 2m$, so $\{i,j\}$ $m$-blocking forces $\mathrm{rank}_i(j), \mathrm{rank}_j(i) \le m$, i.e., $\{i,j\} \in E(G)$; such an edge blocks iff neither endpoint is content w.r.t.\ the other. For $\{u,j\}$ ($j \neq u$), $u$'s inequality $\mathrm{rank}_u(j) + m \le 2m+1$ holds iff $\mathrm{rank}_u(j) \le m+1$. If $\mathrm{rank}_u(j) \le m$ and $u \notin T_j$, then $\mathrm{rank}_j(u) + m \ge 2m+1 > \rho_j$; if $u \in T_j$, then $\{u,j\} \in E(G)$ and non-blocking is exactly $j$ being content w.r.t.\ $u$, included in~(a). If $\mathrm{rank}_u(j) = m+1$, then $j = x_u$ and non-blocking is exactly~(b).
\end{proof}

\section{Reduction to the Residual Case}\label{sec:residual}

Let $M_0$ be any maximal matching of $G$, and let $W := V \setminus V(M_0)$ ($|W|$ is odd). Leaving any $u \in W$ single and pairing $W \setminus \{u\}$ always satisfies Lemma~\ref{lem:crit}(a), since every $G$-edge not in $M$ has an endpoint $i \in V(M_0)$ with $\rho_i \le m$.

Condition~(b) can fail only when $x_u \in W$, $r_u := \mathrm{rank}_{x_u}(u) \le m$, and $x_u$'s partner $v \in W \setminus \{u, x_u\}$ satisfies $\mathrm{rank}_{x_u}(v) \ge r_u + m$. Note that if $|W| = 1$, then $x_u \notin W$ (since $\mathrm{rank}_u(u) = N \neq m+1$), so $M_0$ is $m$-stable.

Suppose no choice of $u \in W$ and pairing of $W \setminus \{u\}$ yields an $m$-stable matching. Then for every $u \in W$, $x_u \in W \setminus \{u\}$, $r_u \le m$, and every $v \in W \setminus \{u, x_u\}$ has $\mathrm{rank}_{x_u}(v) \ge r_u + m$. Hence $|W| \ge 3$, and $x$ has no $2$-cycle on $W$ (if $x_u = w$ and $x_w = u$, then $r_w = \mathrm{rank}_u(w) \le m$, contradicting $\mathrm{rank}_u(x_u) = m+1$). Therefore $x_{x_u} \in W \setminus \{u, x_u\}$, giving $m+1 = \mathrm{rank}_{x_u}(x_{x_u}) \ge r_u + m$, so $r_u = 1$, i.e., $f(x_u) = u$. Thus $W$ has \emph{structure~F}:
\begin{itemize}
  \item $u \mapsto x_u$ permutes $W$ (injective since $f(x_u) = u$) with no $1$- or $2$-cycles;
  \item $f(x_u) = u$ for all $u \in W$; and
  \item every $w \in W$ ranks every $v \in W \setminus \{f(w), w\}$ at position $\ge m+1$.
\end{itemize}

\paragraph{Exchange step.} Suppose $W$ has structure~F and some $w \in W$ has a $G$-neighbor $y \in V(M_0)$ with $z = M_0(y)$. Form $M_1 = M_0 - yz + yw$ and extend $M_1$ to a maximal matching $M_1'$ of $G$. Because $W$ is $G$-independent and $z$ is the only newly uncovered vertex outside $W$, at most one edge $zu$ ($u \in W \setminus \{w\}$) is added. Set $u_0 := x^{-1}(w)$ and $u_1 := x^{-1}(u_0)$; since $x$ has no $1$- or $2$-cycles, $w, u_0, u_1$ are pairwise distinct. The uncovered set $W_1$ of $M_1'$ is odd (hence nonempty). If no edge is added or $zu$ is added with $u \neq u_0$, then $u_0 \in W_1$ has $x_{u_0} = w \notin W_1$; if $zu_0$ is added, then $u_1 \in W_1$ has $x_{u_1} = u_0 \notin W_1$. Hence $W_1$ lacks structure~F, and $M_1'$ extends to an $m$-stable matching.

\paragraph{Residual structure.} Assume henceforth that no maximal matching of $G$ extends to an $m$-stable matching. Then every maximal matching of $G$ leaves uncovered precisely the set $I$ of $G$-isolated vertices, $I$ has structure~F ($|I| \ge 3$ odd), and every maximal matching of $G[J]$ ($J := V \setminus I$) is perfect:
\begin{enumerate}
  \item[(R1)] For every $w \in I$, $T_w \cap I = \{f(w)\}$, so $|T_w \cap J| = m-1$ and every $y \in T_w \cap J$ satisfies $\mathrm{rank}_y(w) \ge m+1$. Thus $n := |J| = 2m+1 - |I|$ satisfies $m-1 \le n \le 2m-2$ (in particular $n \ge 2$).
  \item[(R2)] $G[J]$ is \emph{randomly matchable}: for every edge $uv$ of a component $C$ of $G[J]$, $C - u - v$ admits a perfect $G$-matching. By Sumner's theorem~\cite{Sum79}, every connected component of $G[J]$ is isomorphic to $K_{2r}$ ($r \ge 1$) or $K_{t,t}$ ($t \ge 2$). Let $J_2$ be the union of $K_2$ components ($b := |J_2|$) and $J_+ := J \setminus J_2$ ($a := |J_+|$); every vertex in $J_+$ has $\deg_G \ge 2$.
\end{enumerate}

For each $y \in J$, let $C_y$ be its component in $G[J]$, and define
\begin{align*}
d_y &:= \max_{u \in N(y)} \mathrm{rank}_y(u) \le m, &
A_y &:= \{z \neq y : \mathrm{rank}_y(z) \le d_y + m - 1\}, \\
\varepsilon_y &:= |A_y \cap I|, &
\beta_y &:= \min_{w \in I} \mathrm{rank}_y(w).
\end{align*}
Since $d_y + m - 1 \le 2m - 1 < N$, we have $|A_y| = d_y + m - 1$.

\section{Augmenting Constructions and Property $(\ast)$}

\begin{lemma}[Lemma B]\label{lem:B}
Let $e, e' \in J$ lie in distinct components $C, C'$ of $G[J]$, and let $s \in N(e)$, $t \in N(e')$ satisfy
\[
\mathrm{(i)}\ \mathrm{rank}_e(e') \le \mathrm{rank}_e(s) + m - 1, \qquad
\mathrm{(ii)}\ \mathrm{rank}_{e'}(e) \le \mathrm{rank}_{e'}(t) + m - 1.
\]
For any $w \in I$, form a matching $M$ by leaving $s$ single, matching $e\text{--}e'$ and $t\text{--}w$, taking perfect $G$-matchings of $C-s-e$, $C'-e'-t$, and all other components of $G[J]$, and pairing $I \setminus \{w\}$ so that if $x_s \in I \setminus \{w\}$ then $x_s$ is paired with $f(x_s)$ (when $f(x_s) \neq w$) or $x(x_s)$ (when $f(x_s) = w$). Then $M$ satisfies Lemma~\ref{lem:crit}(a), and~(b) holds unless $x_s \in \{e', t, w\}$. In particular, if $x_s \in I$, choosing $w \in I \setminus \{x_s, f(x_s)\}$ makes $M$ $m$-stable.
\end{lemma}

\begin{proof}
Since $C \neq C'$ and $I$ is $G$-isolated, the only $G$-edges with no $G$-matched endpoint are $\{s,e\}$ (covered by $e$ via~(i)) and $\{e',t\}$ (covered by $e'$ via~(ii)), establishing~(a). For~(b): $x_s \neq e$ since $s \in N(e) \implies \mathrm{rank}_s(e) \le m$; if $x_s$ is $G$-matched then $\rho_{x_s} \le m$; if $x_s \in I \setminus \{w\}$, then either $\rho_{x_s} = 1$, or $f(x_s) = w$ and $x_s$ is paired with $x(x_s) \in I \setminus \{w, x_s\}$ (valid since $x$ has no $1$- or $2$-cycles on $I$, leaving an even set of size $|I|-3$), giving $\rho_{x_s} = m+1 \le \mathrm{rank}_{x_s}(s) + m - 1$ as $s \in J$ implies $s \neq w = f(x_s)$ and hence $\mathrm{rank}_{x_s}(s) \ge 2$.
\end{proof}

\begin{lemma}[Lemma RB]\label{lem:RB}
Let $uy \in E(G)$ and $w \in I$ satisfy:
\begin{itemize}
  \item[\emph{(C1)}] $\mathrm{rank}_y(w) \le \mathrm{rank}_y(u) + m - 1$, and
  \item[\emph{(C2)}] $x_u \neq w$, or $\mathrm{rank}_w(u) \ge m+1$, or $\mathrm{rank}_w(y) \le \mathrm{rank}_w(u) + m - 1$.
\end{itemize}
Then leaving $u$ single, matching $y\text{--}w$, taking a perfect $G$-matching of $J-u-y$, and pairing $I \setminus \{w\}$ as in Lemma~\ref{lem:B} gives an $m$-stable matching.
\end{lemma}

\begin{proof}
The unique $G$-edge without a $G$-matched endpoint is $\{u,y\}$, covered by $y$ via~(C1). Condition~(b) holds for $x_u = w$ by~(C2), and otherwise as in Lemma~\ref{lem:B}.
\end{proof}

Assume from now on that no triple $(u,y,w)$ satisfies Lemma~\ref{lem:RB}.

\begin{proposition}[Property $(\ast)$]\label{prop:star}
If $y \in J$ and $w \in I \cap A_y$, then $C_y = \{u,y\} \cong K_2$, $f(u) = y$, $x_u = w$, $\mathrm{rank}_w(u) \le m$, and $\mathrm{rank}_w(y) \ge \mathrm{rank}_w(u) + m$. In particular, $\varepsilon_y \le 1$ for all $y \in J$, and $\varepsilon_y = 0$ for all $y \in J_+$.
\end{proposition}

\begin{proof}
Let $u \in N(y)$ satisfy $\mathrm{rank}_y(u) = d_y$. Since $w \in A_y$, (C1) holds for $(u,y,w)$, so~(C2) fails: $x_u = w$, $\mathrm{rank}_w(u) \le m$, and $\mathrm{rank}_w(y) \ge \mathrm{rank}_w(u) + m \ge m+1$. For $(y,u,w)$, (C2) holds since $\mathrm{rank}_w(y) \ge m+1$, so~(C1) fails: $m+1 = \mathrm{rank}_u(w) \ge \mathrm{rank}_u(y) + m$, forcing $\mathrm{rank}_u(y) = 1$ ($f(u) = y$). If $u$ had another $G$-neighbor $y' \neq y$, then $\mathrm{rank}_u(y') \ge 2$ gives $\mathrm{rank}_u(w) = m+1 \le \mathrm{rank}_u(y') + m - 1$, so~(C1) holds for $(y',u,w)$; failure of~(C2) yields $\mathrm{rank}_w(u) \ge \mathrm{rank}_w(y') + m \ge m+1$, contradicting $\mathrm{rank}_w(u) \le m$. Hence $\deg_G(u) = 1$, $C_y = \{u,y\}$, and uniqueness of $x_u$ gives $\varepsilon_y \le 1$.
\end{proof}

For $y \in J_2$, denote its unique $G$-neighbor by $y'$ (so $d_y = \mathrm{rank}_y(y')$). Put $K := \{z \in J_2 : \varepsilon_z = 1\}$, $\kappa := |K|$, and for $z \in K$ let $w_z$ be the unique agent in $A_z \cap I$; by Proposition~\ref{prop:star}, $x_{z'} = w_z$ and $z' \in T_{w_z}$.

\section{Good Arcs and Mutual Pairs}

A \emph{good arc} $p \to q$ consists of $p, q \in J$ in distinct components with $q \in T_p$ and $p \in A_q$ (whence $p \notin T_q$, so $\mathrm{rank}_q(p) \ge m+1$). Set $Y_q(p) := \{y \in N(q) : \mathrm{rank}_q(y) \ge \mathrm{rank}_q(p) - m + 1\} \neq \emptyset$.

\begin{lemma}[Lemma GA]\label{lem:GA}
Let $p \to q$ be a good arc.
\begin{enumerate}
  \item For $y \in Y_q(p)$ and $u \in N(p)$, Lemma~\ref{lem:B} with $(e,e',s,t) = (q,p,y,u)$ applies and fails only if $x_y = u$, $y \in T_u$, and $\beta_u \ge \mathrm{rank}_u(y) + m$ (choosing $w \in I$ minimizing $\mathrm{rank}_u(w)$ when $x_y = u \notin I$).
  \item For $u \in N(p)$ and $y \in Y_q(p)$, Lemma~\ref{lem:B} with $(e,e',s,t) = (p,q,u,y)$ applies and fails only if ($x_u = q$, $u \in T_q$, $\mathrm{rank}_q(p) \ge \mathrm{rank}_q(u)+m$) or ($x_u = y$, $u \in T_y$, $\beta_y \ge \mathrm{rank}_y(u)+m$).
\end{enumerate}
Consequently, if $\deg_G(p) \ge 2$, choosing $u \in N(p) \setminus \{x_y\}$ in~(1) yields an $m$-stable matching. If $\deg_G(p) = 1$ ($N(p) = \{u\}$) and both constructions fail for all $y \in Y_q(p)$, then:
\begin{itemize}
  \item[\emph{(G1)}] $\forall y \in Y_q(p)$: $x_y = u$, $y \in T_u$, and $\beta_u \ge \mathrm{rank}_u(y) + m$;
  \item[\emph{(G2)}] $x_u = q$, $u \in T_q$, and $\mathrm{rank}_q(u) \le \mathrm{rank}_q(p) - m$.
\end{itemize}
\end{lemma}

\begin{proof}
Conditions~(i)--(ii) of Lemma~\ref{lem:B} hold by $q \in T_p$ and $y \in Y_q(p)$. In~(1), $x_y = p$ is harmless since $\rho_p = \mathrm{rank}_p(q) \le m$. When $N(p) = \{u\}$ and~(G1) holds, every $y \in Y_q(p)$ satisfies $y \in T_u$, so $\mathrm{rank}_u(y) \le m < m+1 = \mathrm{rank}_u(x_u)$, excluding $x_u = y$ in~(2) and forcing~(G2).
\end{proof}

\begin{lemma}[Lemma K2J]\label{lem:K2J}
Any good arc $p \to q$ with $q \in J_+$ yields an $m$-stable matching.
\end{lemma}

\begin{proof}
If $p \in J_+$, Lemma~\ref{lem:GA} applies. Let $p \in J_2$, $u = p'$, and suppose all constructions fail, so~(G1)--(G2) hold for $p \to q$. By~(G2), $u \in T_q$ and $\mathrm{rank}_u(q) = m+1$. If $d_u \ge 2$, then $m+1 \le d_u + m - 1$, making $q \to u$ a good arc with tail $q \in J_+$, which succeeds by Lemma~\ref{lem:GA}. Thus $d_u = 1$ ($f(u) = p$).

Take $y \in Y_q(p) \subseteq C_q \subseteq J_+$, so $\deg_G(y) \ge 2$ and $d_y \ge 2$. By~(G1), $y \in T_u$ and $x_y = u$, so $\mathrm{rank}_y(u) = m+1 \le d_y + m - 1$; hence $u \to y$ is a good arc with tail $u \in J_2$ (partner $p$), so~(G1)--(G2) hold for $u \to y$. Now~(G2) for $u \to y$ gives $x_p = y$ and $\mathrm{rank}_y(p) \le \mathrm{rank}_y(u) - m = 1$, i.e., $f(y) = p \notin N(y)$. Therefore $Y_y(u) = N(y)$, and~(G1) for $u \to y$ gives $x_{y_2} = p$ for every $y_2 \in N(y)$. Uniqueness of $x_p$ forces $Y_q(p) = \{y\}$. Since $q \in N(y)$, we have $x_q = p$, so $\mathrm{rank}_q(p) = m+1$ and $Y_q(p) = N(q) \setminus \{f(q)\}$. However, (G2) for $p \to q$ gives $\mathrm{rank}_q(u) \le 1$, i.e., $f(q) = u \notin N(q)$, whence $|Y_q(p)| = |N(q)| = \deg_G(q) \ge 2$, contradicting $Y_q(p) = \{y\}$.
\end{proof}

\begin{corollary}\label{cor:S}
If Lemma~\ref{lem:B} applies with single $s \in J_+$, then an $m$-stable matching exists.
\end{corollary}

\begin{proof}
Choosing $w \notin \{x_s, f(x_s)\}$ when $x_s \in I$, failure of~(b) forces $x_s = o \in C' \neq C_s$ with $s \in T_o$. Since $s \in J_+$ has $d_s \ge 2$, $\mathrm{rank}_s(o) = m+1 \le d_s + m - 1$, so $o \to s$ is a good arc into $J_+$ and Lemma~\ref{lem:K2J} applies.
\end{proof}

A \emph{mutual pair} is an unordered pair $\{e, e'\} \subseteq J$ in distinct components with $e' \in A_e$ and $e \in A_{e'}$. Setting $U := \{s \in N(e) : \mathrm{rank}_e(s) \ge \mathrm{rank}_e(e') - m + 1\} \neq \emptyset$ and $Y := \{t \in N(e') : \mathrm{rank}_{e'}(t) \ge \mathrm{rank}_{e'}(e) - m + 1\} \neq \emptyset$, Lemma~\ref{lem:B} applies for all $s \in U, t \in Y$.

\begin{lemma}[Lemmas MJ and D]\label{lem:MD}
Assume no construction above yields an $m$-stable matching. Then:
\begin{enumerate}
  \item Every mutual pair lies in $J_2$ (otherwise $U \subseteq J_+$ and Corollary~\ref{cor:S} applies).
  \item If $\{e, e'\} \subseteq J_2$ is a mutual pair with partner $s = e'$, then $x_s \in C_{e'}$ and $e \notin K$ (since $e \in K$ gives $x_s = w_e \in I$ by Proposition~\ref{prop:star}, whereupon Lemma~\ref{lem:B} succeeds). In particular, every vertex lies in at most two mutual pairs.
\end{enumerate}
\end{lemma}

\section{Counting Bound}\label{sec:counting}

Assume all constructions in Sections~2--4 fail, and let $L$ be the number of mutual pairs (all in $J_2$ by Lemma~\ref{lem:MD}). For $y \in J$, set
\[
\alpha_y := |A_y \cap (J \setminus C_y)| = d_y + m - 1 - \varepsilon_y - |A_y \cap C_y| \ge d_y + m - |C_y| - \varepsilon_y,
\]
since $|A_y \cap C_y| \le |C_y| - 1$ (with equality $= 1$ on $J_2$). Let $P := \frac{1}{2}\sum_C |C|(n - |C|)$ be the number of unordered pairs in $J$ across distinct components. Each such pair contributes $2$ to $\sum_{y \in J} \alpha_y$ if it is mutual and at most $1$ otherwise, so
\[
L \ge \sum_{y \in J} \alpha_y - P \ge \sum_C \left( \sum_{y \in C}(d_y + m - |C| - \varepsilon_y) - \frac{|C|(n - |C|)}{2} \right).
\]
Using $\varepsilon_y = 0$ on $J_+$ and $n \le 2m - 2$ (so $2m - n \ge 2$), each component $C$ contributes:
\begin{itemize}
  \item $C \cong K_2$: $\sum_{y \in C}(d_y - \varepsilon_y) + (2m - 2 - n)$;
  \item $C \cong K_{2r}$ ($r \ge 2$, $d_y \ge 2r-1$): $\ge r(2m - 2 - n) + 2r^2 \ge 2r = |C|$;
  \item $C \cong K_{t,t}$ ($t \ge 2$, $d_y \ge t$): $\ge t(2m - n) \ge 2t = |C|$.
\end{itemize}
Summing over all components yields
\begin{equation}\label{eq:lower}
L \ge a + \sum_{z \in J_2}(d_z - 1) + (b - \kappa) + \frac{b}{2}(2m - 2 - n).
\end{equation}
On the other hand, let $D := \{e' : e \in J_2 \text{ lies in a mutual pair}\} \subseteq J_2$. By Lemma~\ref{lem:MD}(2), each endpoint $e$ lies in at most $2$ mutual pairs and $e \notin K$ (so $e' \notin K' := \{z' : z \in K\}$). Since $e \mapsto e'$ is a bijection on $J_2$, we get $2L \le 2|D|$ and $D \subseteq J_2 \setminus K'$, hence
\begin{equation}\label{eq:upper}
L \le |D| \le b - \kappa.
\end{equation}
Combining~\eqref{eq:lower} and~\eqref{eq:upper} gives $a + \sum_{z \in J_2}(d_z - 1) + \frac{b}{2}(2m - 2 - n) \le 0$. Since every term is nonnegative and $n \ge m-1 > 0$, we conclude:
\[
a = 0 \ (J = J_2, \ b = n > 0), \qquad d_z = 1 \ (\forall z \in J), \qquad n = 2m - 2 \ (|I| = 3).
\]
Because $d_z = 1$, we have $A_z = T_z$ for all $z \in J$, so no mutual pair can exist between distinct components: $L = 0$. Then~\eqref{eq:lower} forces $\kappa = b$, so every $z \in J$ has a unique $w_z \in T_z \cap I$, and $\alpha_z = m - 2$. Finally, $\sum_{z \in J} \alpha_z = (2m-2)(m-2) = P$, so every pair of non-partners in $J$ contributes \emph{exactly}~$1$ (forming a tournament on non-partners). Write $\mathrm{out}(z) := T_z \cap (J \setminus \{z, z'\})$; then $|\mathrm{out}(z)| = m - 2$.

\section{The All-$K_2$ Case ($m \ge 3$)}\label{sec:allk2}

For each $z \in J$, set $c(z) := x_z \in I$. Applying Proposition~\ref{prop:star} to $y = z'$ and $w = w_{z'} \in T_{z'} \cap I$ gives $c(z) = w_{z'}$, $z \in T_{c(z)}$, and $\mathrm{rank}_{c(z)}(z') \ge \mathrm{rank}_{c(z)}(z) + m > m$. Hence partners have distinct colors: if $c(z) = c(z') = q$, then both $z, z' \in T_q$, contradicting $\mathrm{rank}_q(z') > m$.

Fix any $z \in J$ and $a \in \mathrm{out}(z)$ (so $z, z', a, a'$ are four distinct agents in $J$). Form $M$ by leaving $z'$ single, matching $z\text{--}a$ and $a'\text{--}c(a)$ (note $c(a) = w_{a'} \in T_{a'}$), matching every remaining $K_2$ component as a $G$-edge, and pairing the two elements of $I \setminus \{c(a)\}$.
\begin{itemize}
  \item \emph{Lemma~\ref{lem:crit}(a):} the only $G$-edges without a $G$-matched endpoint are $\{z, z'\}$ ($z$ has $\rho_z = \mathrm{rank}_z(a) \le m$) and $\{a, a'\}$ ($a'$ has $\rho_{a'} = \mathrm{rank}_{a'}(w_{a'}) \le m = \mathrm{rank}_{a'}(a) + m - 1$).
  \item \emph{Lemma~\ref{lem:crit}(b):} the single is $z'$ with $x_{z'} = c(z') \in I$. If $c(z') \neq c(a)$, then $c(z')$ is paired with the unique other element of $I \setminus \{c(a)\}$. Since $|I| = 3$ and $x$ is a $3$-cycle on $I$, this partner is either $f(c(z'))$ ($\rho_{c(z')} = 1$) or $x(c(z'))$ ($\rho_{c(z')} = m+1 \le \mathrm{rank}_{c(z')}(z') + m - 1$, as $f(c(z')) = c(a) \in I \neq z'$ implies $\mathrm{rank}_{c(z')}(z') \ge 2$).
\end{itemize}

Therefore an $m$-stable matching exists unless $c(a) = c(z')$ for every $z \in J$ and every $a \in \mathrm{out}(z)$. Assume for contradiction that this holds. Because $m \ge 3$, $|\mathrm{out}(z)| = m - 2 \ge 1$, so pick $z \in J$ and $a \in \mathrm{out}(z)$, and consider the non-partner pair $\{z, a'\}$:
\begin{itemize}
  \item If $a' \in \mathrm{out}(z)$, then $c(a') = c(z') = c(a)$, contradicting $c(a') \neq c(a)$.
  \item Otherwise, by the tournament property, $z \in \mathrm{out}(a')$, so $c(z) = c((a')') = c(a) = c(z')$, contradicting $c(z) \neq c(z')$.
\end{itemize}
This contradiction completes the proof of Theorem~\ref{thm:main}. \qed

\begin{thebibliography}{9}
\bibitem{FM25}
A.~Fanelli and L.~Moscardelli, \emph{Approximately Stable Matching}, in Proc.\ ECAI 2025, Frontiers in Artificial Intelligence and Applications, IOS Press, 2025, \url{https://doi.org/10.3233/FAIA251234}.

\bibitem{Sum79}
D.~P.~Sumner, \emph{Randomly matchable graphs}, J.\ Graph Theory \textbf{3} (1979), 183--186.
\end{thebibliography}

\end{document}
